\section{Algebraic closure of the zero Fourier mode associated with CT108 periodicity in the electron fibre}
\label{sec:clausura-nulo-material-electron-rev3}
\index{electron!null--material closure}
\index{CT108!null mode and electron}

Planck's self-dual identity fixes the origin of the positive line of
masses; the electronic genealogy fixes its first complete material
displacement. This section brings together, without repeating their reference proofs,
the circular realization of CT108, the basal composition of the electron, and the
bidirectional holonomy of its route.

Let
\[
 \mathcal H_{108}^{\rm reg}
 :=\ell^2(\mathbb Z/108\mathbb Z;\mathbb R),
 \qquad
 S_{108}^{\rm reg}\delta_\theta=\delta_{\theta+1},
\]
the real regular representation of the cyclic shift of CT108, and let
\[
 n_0:=\frac1{\sqrt{108}}\sum_{\theta=0}^{107}\delta_\theta,
 \qquad
 \mathcal N_{108}:=\ker(S_{108}^{\rm reg}-I)=\mathbb R n_0.
\]
The normalized zero Fourier mode is thereby defined, not merely named.
Let \(\mathcal G_e\) be the space of enriched electronic genealogies of
\cref{prop:planck-ct108-realizacion-centro-electron-rev3}, and fix the
linear submodule
\(\ell_{\Gamma_e}:=\mathbb R g_{\Gamma_e}
\subset\mathbb R[\mathcal G_e]\) generated by the central route in the
free linearization of that space. Likewise, if
\(F_e^+=\mathbb R_{>0}I\) is the positive ray already defined, let
\(\mathcal F_e:=\operatorname{span}_{\mathbb R}(F_e^+)=\mathbb R I\) be its
linear span, and let
\(\mathcal L_M\) be the mass line defined in
\cref{sec:ley-general-masa-accion-autonoma}. The declared bases determine
the rank-one linear morphism
\begin{equation}
 \mathfrak C_e:
 \mathcal N_{108}\otimes\ell_{\Gamma_e}
 \longrightarrow \mathcal F_e\otimes\mathcal L_M,
 \qquad
 \mathfrak C_e(n_0\otimes g_{\Gamma_e})
 =\mathfrak c_e\,I\otimes\mathcal U_M,
 \label{eq:morfismo-clausura-nulo-material-electron-rev3}
\end{equation}
where the positive coordinate is
\begin{equation}
 \boxed{
 \mathfrak c_e=
 \frac{\sqrt3}{4}
 \left[
  \exp\!\left(\frac{\varphi_{\rm HMT}}{\pi_{\rm HMT}^{2}}\right)
  -22\alpha_{\rm HMT}^{3}
 \right]
 (1+15\Delta_4)
 R_{\rm act}^{\,80-54\alpha_{\rm HMT}+6\Delta_4}.}
 \label{eq:coordenada-clausura-nulo-material-electron-rev3}
\end{equation}

\begin{theorem}[Null closure--one-range electronic material]
\label{thm:clausura-nulo-material-electron-rev3}
With the primitives and charts already fixed, the morphism \(\mathfrak C_e\) exactly factors the beta electron output. Its four factors preserve, in this order, the sum--product incompatibility of APP; the areal--volumetric orientation together with the boundary deficit; the torsion and retained memory; and the bidirectional holonomy of the electron route. In the chart \(\mathcal U_M=1\,\mathrm{MeV}/c^2\),
\begin{equation}
 \mathfrak C_e(n_0\otimes g_{\Gamma_e})
 =0.5109989506900008086082571648\ldots\,
  I\otimes\mathrm{MeV}/c^2.
 \label{eq:masa-clausura-nulo-material-electron-rev3}
\end{equation}
\end{theorem}

\begin{proof}
The \cref{thm:cierre-electron-holografico-rev6} proves that the first three
factors of
\eqref{eq:coordenada-clausura-nulo-material-electron-rev3} compose the
positive basal scalar \(\mathfrak e_{\rm HMT}\). The
\cref{thm:reduccion-electron-two-way-rev11} obtains, from two independent
readers of the route,
\[
 K_D(\Gamma_e)=K_\Omega(\Gamma_e)=(80,54,6),
 \qquad
 \kappa_e=80-54\alpha_{\rm HMT}+6\Delta_4.
\]
The ratio \(R_{\rm act}>0\) belongs to the unique pre-return/post-return
action line. Consequently, the fourth factor is positive, and
the composition coincides with
\eqref{eq:electron-beta-two-way-rev9} and
\eqref{eq:masa-electron-two-way-rev11}. The dimensional chart acts only
after this factorization.
\end{proof}

\paragraph{Mathematical scope and physical realization.}
The exact result is the declared linear realization of the electron factorization over the typed null mode, genealogy, and fibre. The CT108 null mode does not by itself select the coordinate \(\mathfrak c_e\): that coordinate follows from the four factors proved in their respective residences, and the morphism assembles them into a single rank-one algebraic closure. The bound electromagnetic dynamics retains as its own requirements a stable Hamiltonian, charge, spin \(1/2\), the form factor, and gauge coupling; these objects are not premises of the preceding factorization. The Planck unit occupies the self-dual point of the line, whereas the electron occupies a genealogically displaced section: an interpretation as a classical black hole belongs to another type of realization and does not enter this morphism.
